\subsection{Smooth covers at a crossing}\label{p10-sec:incidence}

Let $f:Y\to C$ be a smooth connected cover, and write $D_f=f^*D$.
We regard covers up to isomorphism over $C$. Let $\mathcal H$ denote
the fixed-target Hurwitz space of the ramification type of $f$.
At a point $(f,t)$ with kernel $K$ and cokernel $P$ the differential
of the family of operators induces a map
\begin{equation}\label{p10-eq:cubic-intrinsic}
 W:T_f\mathcal H\oplus\RR_t\longrightarrow
                  \operatorname{Hom}_{\RR}(K,P).
\end{equation}
It is obtained by differentiating on $K$ and then dividing by the
image of $D_f$, so it is independent of the trivializations. When
$K$ is a line and $W$ is an isomorphism, write $\eta_f$ for its sign.
We orient the source by the complex orientation of the Hurwitz space
and increasing time. On the target we use the full Cauchy--Riemann
determinant orientation, continued from complex-linear operators.
Indeed $\dim_{\RR}P=2\dim_{\CC}\mathcal H+1$ in this case, and
$\det\operatorname{Hom}_{\RR}(K,P)$ differs from
$(\det D_f)^*$ by an even tensor power of $K^*$. This convention
also applies to a double cover.

\begin{proposition}\label{p10-prop:incidence}
Fix endpoints rigid through degree three. For a residual set of smooth
paths between them with fixed symbol and invertible base operator,
every smooth cubic cover with kernel has monodromy $S_3$, trivial deck
group, and only simple ramification over distinct branch values. Its
kernel is a real line, the graph of a generator is embedded, and
\eqref{p10-eq:cubic-intrinsic} is an isomorphism. At that parameter every
other smooth connected cubic cover and every smooth connected double
cover have injective pullback, in every source genus.
\end{proposition}

\begin{proof}
We first vary the operator and a nonzero kernel vector $u$ on one
cubic cover. On an unramified disk write $u_1,u_2,u_3$ for its three
values. Their sum is single-valued on the punctured target and
solves $D$ there. It is bounded at the branch points and hence
extends across them by removal of singularities. Base invertibility
gives $u_1+u_2+u_3=0$. On the connected Galois closure these three
sections span the real sum-zero representation of the monodromy
group. In fact the map from the permutation representation sending
the coordinate vectors to the three conjugate sections kills the
invariant vector, and the remaining real plane is irreducible for
either transitive group $C_3$ or $S_3$. In particular each pair of
sheet values is real-linearly independent on an open dense subset.
A relation $u_i=a u_j$ on an open set with $u_j\ne0$ and $a$ real
gives $(\bar\partial a)u_j=0$, so $a$ is constant. Unique
continuation on the Galois closure contradicts the representation
just computed.

Consider the universal equation $D_fu=0$. A covector $p$ annihilating
its differential in the section and operator directions solves the
homogeneous adjoint equation. On an unramified disk use the base
coordinate to write $p_i$ for its sheet values, with the density
in the real duality pairing expressed in that coordinate.
An arbitrary real zeroth-order variation $A$ of $D$ supported there
gives the variations $(Au_1,Au_2,Au_3)$. Where the first two vectors
are real-linearly independent, these range over all triples of form
values with sum zero. Thus $p_1=p_2=p_3$ on such a disk. Testing the
section variation $f^*w$ gives $\langle p,f^*(Dw)\rangle=0$ for
all smooth $w$. Every smooth base form equals $Dw$ by base
invertibility, so $\sum_i p_i=0$ on every unramified disk. Hence
$p$ vanishes there and identically by unique continuation. The
differential is surjective since its range contains the closed
finite-codimensional image of $D_f$.

Suppose the cover has $b$ simple and $q$ totally ramified values.
These exhaust the cubic profiles. In a chart with ordered branch
values, the $b+q$ distinct values are local complex coordinates.
Indeed the finite covering of their complement is locally fixed by
monodromy and has a unique completion by the prescribed power maps.
Riemann--Roch and Riemann--Hurwitz give
$\operatorname{ind}_{\RR}D_f=-2(b+2q)$. Quotienting the vector by
nonzero real scaling removes one kernel direction. The projectivized
incidence thus projects to the space of operators with index
\[
                    2(b+q)-2(b+2q)-1=-1-2q .
\]
Relative path transversality excludes $q>0$ and leaves a regular
zero-dimensional incidence for $q=0$. A kernel of dimension greater
than one would give a positive-dimensional projective fiber. A cyclic
cubic has no invariant kernel vector, by descent to the invertible
base, and every nontrivial irreducible real representation of $C_3$
has dimension two. Thus the monodromy is $S_3$. A point stabilizer
in $S_3$ is self-normalizing, so the deck group is trivial. At the
remaining covers, projection of the regular linearization to $P$ gives
the surjectivity of $W$. Both sides have dimension $2b+1$.

We give the marked transversality argument because separation of two
cubic covers will require a condition on their complex sheet spans.
At distinct points of a fiber evaluate the difference of their values.
At a ramification point take the complex-linear first derivative of
the section in a frame pulled back from $C$. Its antiholomorphic
derivative already vanishes at a solution since $df=0$ there.
Finally, over a fixed regular base point, evaluate the pair
\[
                     (u_1-u_3,u_2-u_3)\in N_x^2.
\]
Each of these linear maps $J$ on sections kills all base pullbacks.
To check surjectivity for the equation augmented by $J$, hold the
marks fixed and let $(p,\gamma)$ annihilate the section and operator
variations. Off the marks $p$ solves the homogeneous adjoint equation.
The same test with base pullbacks kills its trace, and the localized
variations $A$ force $p$ to vanish on an unmarked disk. By unique
continuation $p$ is supported on the marks. A distribution supported
at these finitely many points has finite order. At each mark we may
prescribe the value and the holomorphic first derivative of a test
section $w$ while arranging that $D_fw$ vanish to any fixed jet order.
In a local expression $D_f=\bar\partial+B$ the Taylor coefficients
with positive antiholomorphic degree are determined successively,
leaving the holomorphic coefficients free. Multiplying the resulting
polynomials by cutoffs gives the required global sections. Now
$\langle p,D_fw\rangle+\gamma(Jw)=0$ forces $\gamma=0$. Elliptic
regularity for the remaining adjoint equation gives $p=0$.

The coincidence of two values has real codimension four. The base
point in an unramified fiber varies in two real dimensions, and a
mark over a branch value, including the pair consisting of the
ramified and the unramified preimage, is discrete over its cover.
Vanishing of the derivative at a ramification point likewise imposes
four real conditions at a discrete mark. A generic path avoids all
three defects, so the graph is an injective immersion of a compact
curve. For the last evaluation the pair has complex rank one in real
codimension two and rank zero in codimension eight. Taking a countable
dense set of fixed base points shows that at every crossing there
is an unramified disk on which the two differences, and hence
$u_1,u_2$, are complex-linearly independent. The rank condition is
unchanged by relabeling the sheets.

It remains to exclude a second cover at the same time. Use separate
Hurwitz and section variables and a common variation $A$ for the
two equations. An annihilating pair of adjoints solves the homogeneous
adjoint equation on each cover and has zero trace on each by the
base-pullback test. For a nonzero cubic adjoint the three values on
an unramified disk span a real plane on an open dense set. Indeed,
view the adjoint there as a dual-bundle-valued differential in the
base coordinate. On the connected regular cover of the complement of
the branch set the trace-free conjugates form the same irreducible
plane as above. Real proportionality on a disk forces a constant
ratio by their common adjoint equation. The argument needs no
extension of the adjoint divided by the base differential across a
branch point.

The tensor pairing with $A$ on that disk is
$\sum_i u_i\otimes p_i$. Substituting the trace identities gives
coefficient matrix
$\left(\begin{smallmatrix}2&1\\1&2\end{smallmatrix}\right)$ in the
first two values of each triple. Thus, when both pairs are
real-linearly independent, the tensor has real rank two and image
$\operatorname{span}_{\RR}(u_1,u_2)$. For a double cover the
corresponding tensor is $2u_1\otimes p_1$, of rank one on an open
dense set for nonzero adjoint. Tensors of these different ranks
cannot cancel. This proves surjectivity for a cubic and a double
cover with marked kernel lines.

For two cubics restrict to the open locus where the first kernel
has a complex-independent pair on some disk. If both adjoints
are nonzero, cancellation for all $A$ forces the two real sheet
planes to agree on a smaller disk on which all the rank conditions
hold. Express any kernel sheet of the second cover as $a u_1+b u_2$
there, with $a,b$ real. Its equation gives
$(\bar\partial a)u_1+(\bar\partial b)u_2=0$. Complex independence
makes $a,b$ constant. On the connected regular cover of the common
punctured target defined by the intersection of the two monodromy
kernels, unique continuation extends all these relations. The two
planes of conjugate sections coincide as representations of the
punctured target's fundamental group. The sum-zero representation
of a transitive cubic monodromy group is faithful, so the monodromy
kernels agree. In their common quotient, all index-three subgroups
are conjugate, both for $S_3$ and for $C_3$. The two covers are
isomorphic on the punctured curve and then on its unique normal
completion. This excludes cancellation for distinct covers. If
only one adjoint were nonzero, its tensor would already have
positive rank. The paired universal equation is therefore
transverse off the diagonal. Its projection index is the sum of the
two single-cover indices, at most $-2$. Here the double-cover index
is $-1$ by the same branch count and Riemann--Hurwitz. A generic
path avoids these pairs.

All uses of transversality may be made on sheet-labelled Hurwitz
charts, with Sobolev order large enough for moving first jets and a
higher differentiability class for coefficient variations. Scaling
is a free kernel direction before taking projective lines. The
equations augmented by evaluations have closed finite-codimensional
range in the section variables. Consequently smooth coefficient
variations suffice for all the surjections above. More explicitly,
exhaust each marked chart by compact subsets away from the relevant
diagonals and normalize kernel vectors in $L^2$. On each compact
subset elliptic estimates give precompactness of the vectors for
smoothly converging paths. Transversality there is open and can be
achieved by finite-dimensional smooth perturbations supported away
from the endpoints, by the parametric transversality theorem. For the
open locus in the paired calculation one also exhausts the locus
where a specified complex minor is nonzero. Countably many such
charts, profiles and genera suffice. This yields the asserted residual
set in the smooth topology. For the analogous double-cover marked
incidence see \cite[Lemmas~5.11 and~5.15]{bai}.
\end{proof}

\begin{proposition}\label{p10-prop:smoothjump}
Suppose a path as in Proposition~\ref{p10-prop:incidence} meets the wall
at a smooth cubic cover $f:Y\to C$ of genus $h$, at time zero. There
is a punctured interval of operators rigid through degree three. For
any sufficiently small $t_-<0<t_+$ the series on the original curve
satisfy
\[
 L_3(D_{t_+};v)-L_3(D_{t_-};v)=2\eta_f x^{h-1},
 \qquad
 L_d(D_{t_+};v)=L_d(D_{t_-};v)\quad(d=1,2).
\]
At a simple double crossing of genus $h$, the corresponding
degree-two change is $\eta_f x^{h-1}$.
\end{proposition}

\begin{proof}
Fix a generator $\xi$ of $K$ at the cubic cover. For a simply
branched cover the quotient of $TY\to f^*TC$ is torsion of length
$b$, the number of simple branch values. Its $H^0$ is the Hurwitz
tangent and its $H^1$ vanishes. Thus the tangential deformation
problem is regular and contributes the complex orientation.

We realize the path in a neighborhood of the zero section by
prescribing the graph equation over a holomorphic projection
$\varphi:Y'\to C$ to be
\begin{equation}\label{p10-eq:cubic-graphequation}
 \varphi^*D_t u+\varphi^*B(u)=0 .
\end{equation}
Here $B$ is a smooth fiberwise real quadratic with values in
$\Omega^{0,1}_C\otimes N$ and $\varphi^*$ pulls back its form argument.
In a connection splitting, the horizontal-to-vertical block of the
almost complex structure prescribes precisely the zeroth-order term
in this equation, with standard complex structure on the fibers
\cite[Lemma~5.9]{bai}.
Cut off this block outside a disk to extend the structures to
$\PP(N\oplus\CC)$. A fiber Fubini--Study form extended using a
unitary connection, plus a sufficiently large multiple of a positive
base form, tames the compact family. The extension can be made for
all fiber dilations used below on a fixed disk of radius greater than
one, by applying the cutoff outside that disk.

Choose a nonzero adjoint $p$ annihilating
$W(T_f\mathcal H)(\xi)\subset P$. On an unramified disk where one
of its values $p_1$ is nonzero and the three $\xi_i$ lie on
distinct real lines, choose smooth real covectors $\alpha_2,\alpha_3$
with $\alpha_i(\xi_i)=0$ and $\alpha_i(\xi_1)=1$. Multiplying
$\alpha_2(u)\alpha_3(u)$ by a compactly supported vector-valued
form pairing positively with $p_1$ gives $B$ with
$\langle p,f^*B(\xi)\rangle\ne0$. Only the first sheet contributes.
This choice is possible for a cubic because its deck group is trivial.

Use local Hurwitz coordinates $y$ and a real kernel amplitude $z$.
Write $u=z(\xi+w)$ with $w$ in a fixed complement of $K$. After
division by $z$ the image equation determines $w=w(y,t,z)$ by the
implicit function theorem. The remaining equation is $F(y,t,z)=0$
in $P$. Its $(y,t)$ derivative at zero is $W(\cdot)(\xi)$, and
$F_z$ pairs nontrivially with $p$ by the choice of $B$. The terms
from differentiation of $w$ lie in the image of $D_f$. Thus the
solution is $(y(z),t(z))$, where $t(0)=0$ and $t'(0)\ne0$. The
undivided reduced equation is $zF=0$. There is one nonzero graph
on each nearby slice, and it is embedded since dilation by $1/z$
makes it converge smoothly to $(f,\xi)$.

For the fixed-time reduced section $zF$, order its columns by the
complex-oriented $y$ variables followed by $z$. Along the nonzero
branch the first $2b$ columns are $zF_y$, and modulo these the last
is $-zt'(z)F_t$. The identifications of the invertible image
complement and the tangential determinant are the same as those
used in \eqref{p10-eq:cubic-intrinsic}. The normal determinant sign of
the embedded graph is consequently
\[
             -\operatorname{sgn}\bigl(z^{2b+1}t'(z)\bigr)\eta_f
                    =-\operatorname{sgn}(t)\eta_f .
\]
Changing $\xi$ changes the orientations induced on $z$ and on $P$
together. The signed number of nonzero curves therefore changes by
$-2\eta_f$. This is the smooth-domain bifurcation of
\cite[Theorem~4.12(ii)]{bai}, realized here for the prescribed path
of first-order operators.

We next isolate these supports in all normalization genera. Suppose
there are additional simple curves of degree at most three contained
in tubes shrinking to the zero section at parameters tending to zero.
Divide their heights by their fiber maxima $\rho>0$. Their equations
on the unit disk have quadratic term $\rho B$ and converge smoothly
to the linear graph equation. Their image cycles have bounded mass
for a common taming form, since its class on the disk bundle comes
from the base and the projection degrees are bounded. Cycle compactness
gives a limiting holomorphic cycle with convergence of currents and
Hausdorff convergence of supports \cite[Theorem~2.11]{diw}. Integration
of a pulled-back area form preserves projection degree, and the limit
contains a point of height one. Every irreducible component has
positive projection degree, since a compact complex curve in an
affine fiber is constant. The normalization of a component reaching
positive height therefore carries a nonzero pullback kernel. By
Proposition~\ref{p10-prop:incidence} this component is the graph of a
nonzero multiple of $\xi$. Its degree exhausts the available three.
The limit is thus embedded of multiplicity one.

Near such a cycle the simple parametrizations are embeddings and
converge up to reparametrization
\cite[Theorem~2.20 and Proposition~2.21]{diw}. In particular their
genus stabilizes, and elliptic regularity gives smooth convergence.
After scaling back, the covers and sections belong to the Hurwitz
chart just analyzed. This proves the existence of one closed tube
and one interval with no additional simple supports. Shrink the
interval so the nonzero branches lie strictly inside. For the linear
endpoint equation use the same argument with unit-height kernel
graphs for a sequence of nonzero times tending to zero, retaining the
mapping multiplicity in each pushforward cycle. The embedded
multiplicity-one limit again gives convergence of the covers and
normalized sections. The reduced equation for this linear problem is
$F(y,t,0)=0$, whose unique nearby solution is $(0,0)$. Thus on
the punctured interval all smooth pullbacks through degree three are
injective. On a nodal domain every positive normalization then has
zero section and the matched constants on contracted components are
zero as well.

Fix an arithmetic genus $k$ and a positive degree $d\le3$. The
stable maps whose images lie in the closed tube form an open and
closed part of the parametrized moduli in the compact completion.
Indeed all their simple supports lie in the interior, and compactness
preserves containment in the closed tube. Their endpoint virtual
degrees agree by deformation invariance
\cite[\S3.1, (3.4)--(3.5)]{diw}. Maps supported on $C$ contribute
the fixed-target Euler number by
\cite[Theorem~1.2, (1.7) and~(1.10)]{zinger}. Using even one nonzero graph
exhausts degree three, so a map with that support has degree one onto
the embedded graph, with arbitrary contracted components.

For clarity, let $D'$ be the invertible normal operator of one such
embedded curve. A degree-one stable map to the curve has one effective
normalization mapping isomorphically, and each connected contracted
subcurve attaches at just one point. For any deformation of $D'$,
restriction to the effective normalization identifies the kernels
of the base and pullback operators and is a left inverse on their
cokernels. Thus the quotient of the two Fredholm complexes has
zero kernel throughout, and the stabilized quotient cokernels form a
bundle, including over nodal maps \cite[Appendix~B]{bai}.
At an invertible endpoint its orientation differs from the full
normal orientation by the determinant sign of $D'$. The complex
normal degree is $2h-2$, since the ambient first Chern class vanishes
on the disk bundle. Deform to the Dolbeault operator of two acyclic
lines of degree $h-1$. The quotient bundle there has complex
orientation, and its degree-one Euler series is $x^{h-1}$ by
\cite[Theorem~2]{pandhodge}. Consequently the change contributed by
the nonzero graphs is $-2\eta_f x^{h-1}$, including every contracted
configuration at its own fixed arithmetic genus. Conservation gives
the asserted change on $C$ and constancy in lower degrees.

For a double cover the involution negates its kernel generator and
identifies the two nonzero amplitudes in an odd cubic perturbation.
The fixed-target comparison then gives the change $\eta_f x^{h-1}$
in degree two (Theorem~\ref{p9-thm:wall}).
\end{proof}
